

\begin{frame}
  \frametitle{SNM diversion may be detected using reactor kinetics}
    \begin{block}{\Gls{SNM} diversion scenarios considered are \cite{protection2011evaluation, international2022iaea, bathke2012attractive}:}
        \begin{itemize}
            \item One \gls{SQ} of $^{235}$U (75 $kg$) after one year in the \gls{MSRE}
            \begin{itemize}
                \item Removing $^{235}$U will decrease the delayed neutron fraction
            \end{itemize}
            \item One \gls{SQ} of Pu (8 $kg$) after one year in the \gls{MSRE}
            \begin{itemize}
                \item Removing Pu will increase the delayed neutron fraction
            \end{itemize}
        \end{itemize}
    \end{block}
    The \gls{DNP} group parameters will then be used in a negative reactivity step insertion, and the resulting power responses will be compared \cite{wheeler2021signatures}.
    \begin{block}{Each scenario will have four sets of \gls{DNP} group parameters to compare:}
    \begin{itemize}
        \item Conventional parameters with/without diversion
        \item Improved parameters with/without diversion
    \end{itemize}
    \end{block}

    \textbf{The change in fuel composition and the \gls{MSR} phenomena affect the ability to detect \gls{SNM} diversion.}

\end{frame}
